\section{Objectives and hypotheses}
\label{sec:objectives}

The object of the evaluation is the \FS{} solver: a combinatorial simplex whose
basis is represented by a forest on the precedence graph, applied to the
normalised maximum-ratio subproblem and iterated on the residual to build the
canonical sequence of blocks that yields the LP value at any capacity.

Two further backends were implemented for this report and are used as
references, not as contributions: a Dinkelbach iteration whose ratio oracle is a
maximum-closure minimum cut, and a divide-and-conquer parametric decomposition
of the same closure family. All three produce the same certified structures and
are checked by the same verifiers.

The campaigns answer the following questions.

\begin{description}[leftmargin=2.4em,style=nextline]
  \item[E1 Correctness] Does the solver produce certified optimal solutions on
    generic DAGs, with negative profits, ties and degeneracy? Do the three
    backends and an independent LP solver agree?
  \item[E2 Overall performance] How do the reference \FS{} configuration, the
    tuned one and the two combinatorial baselines compare on a frozen test set,
    per family rather than only in aggregate?
  \item[E3 Scalability] How do time and memory grow with $n$ and $m$, and how
    much of the growth is the cost of a pivot rather than the number of pivots?
  \item[E4 Ablation] What does each implementation choice contribute: local
    forest updates, dual canonicalisation, partial pricing, transitive
    reduction, vertex ordering?
  \item[E5 Tuning] Does the search produce a default that transfers from
    training and validation to instances never used to choose it?
  \item[E6 One capacity against many] When does building the whole positive path
    pay for itself against a single solve, once the cost of the queries and of
    materialising the solution is included?
  \item[E7 Degeneracy and numerics] What is the cost of ties, exact and
    near-exact, and of the anti-cycling fallback?
  \item[E8 Applied instances] What happens on the historical PCKP instances and
    on the open-pit instances derived from MineLib?
\end{description}

No question is posed so that a particular answer is the only admissible one. In
particular, the report states plainly where the \FS{} implementation is slower
than the baselines, and section~\ref{sec:discussion} discusses why.
